% Standalone notation bridge for the current workbench spectral modules.
% The root paper uses R for the logarithmic radial height and T=Im(tau).

\section{Canonical radial notation for the current operator modules}
\label{cusp-spec-section}

The earlier sections use \(R\) for the logarithmic cusp height and
\(T=\operatorname{Im}\tau\).  The current workbench operator modules use the
logarithmic height as their independent parameter.  In their standalone
adaptations below that parameter is consistently denoted by \(R\), so no
second meaning is assigned to \(T\).  Thus
\begin{equation}
 t_c=e^{-2\pi R},\qquad s=iR,\qquad R\ge R_0,
 \label{eq:current-radial-ray}
\end{equation}
and we set
\begin{equation}
 \begin{aligned}
  \mu_R&:=\mu(e^{-2\pi R})=\alpha_R+i m_R,\\
  h_R&:=h(e^{-2\pi R})=\eta_R+i H_R,\\
  b_R&:=b(e^{-2\pi R})=\xi_R+i B_R,\\
  c_R&:=\xi_R-\eta_R,\qquad
  L_R:=R+H_R,\qquad q_R:=L_R-B_R,\\
  D_R&:=L_Rq_R+6m_R^2.
 \end{aligned}
 \label{cusp-spec-radial-notation}
\end{equation}
All quantities on the right are real.  In particular
\(L_R=\operatorname{Im}\tau_R\),
\(-q_R=\operatorname{Im}\beta_R\), and \(D_R>0\) is the physical covolume,
whereas the determinant function \(D\) in \eqref{eq:D} equals
\(-D_R/L_R<0\).  Within the imported operator modules below, bare \(m\)
denotes the limiting constant, not the varying function in
\eqref{eq:D}:
\begin{equation}
 m:=\operatorname{Im}\mu(0)=m_0,\qquad
 H:=\operatorname{Im}h(0),\qquad
 B:=\operatorname{Im}b(0)=b_{\mathrm I,0}.
 \label{cusp-spec-limiting-imaginary-parts}
\end{equation}
Holomorphy at \(t_c=0\) gives, without identifying any of these three
constants,
\begin{equation}
 m_R=m+O(e^{-2\pi R}),\qquad
 H_R=H+O(e^{-2\pi R}),\qquad
 B_R=B+O(e^{-2\pi R}).
 \label{eq:current-radial-errors}
\end{equation}

We retain the ordered real coordinates
\begin{equation}
 (z_1,z_2)=(x_1+iy_1,x_2+iy_2)
 \longmapsto(x_1,y_1,x_2,y_2).
 \label{cusp-spec-real-coordinate-order}
\end{equation}
They carry the unrescaled Euclidean metric used throughout the operator
modules.

\begin{proposition}[Exact translation to canonical radial variables]
\label{cusp-spec-exact-dual-lattice}
In the ordered marking
\((\widehat\gamma,\widehat u,\widehat w,\widehat\delta)\), the real period
matrix is
\begin{equation}
 P_R=
 \begin{pmatrix}
  6\alpha_R&\eta_R&1&0\\
  6m_R&L_R&0&0\\
  c_R&\alpha_R&0&1\\
  -q_R&m_R&0&0
 \end{pmatrix}.
 \label{cusp-spec-real-period-matrix}
\end{equation}
Its determinant and positive covolume are
\begin{equation}
 \det P_R=-D_R,\qquad
 \operatorname{covol}(P_R\mathbb Z^4)=D_R>0,
 \qquad
 \operatorname{Vol}(F_R)=D_R,
 \label{cusp-spec-exact-volume}
\end{equation}
where \(F_R=\mathbb R^4/P_R\mathbb Z^4\).  The complete inverse transpose is
\begin{equation}
 P_R^{-\mathsf T}=
 \begin{pmatrix}
 0&0&1&0\\[2pt]
 \dfrac{m_R}{D_R}&\dfrac{q_R}{D_R}&
 \dfrac{-6m_R\alpha_R-q_R\eta_R}{D_R}&
 \dfrac{-m_Rc_R-q_R\alpha_R}{D_R}\\[7pt]
 0&0&0&1\\[2pt]
 -\dfrac{L_R}{D_R}&\dfrac{6m_R}{D_R}&
 \dfrac{6(L_R\alpha_R-m_R\eta_R)}{D_R}&
 \dfrac{L_Rc_R-6m_R\alpha_R}{D_R}
 \end{pmatrix}.
 \label{cusp-spec-full-inverse-transpose}
\end{equation}
For \(k=(r,s,n_1,n_2)^{\mathsf T}\in\mathbb Z^4\), define
\begin{equation}
 \mathcal R_R(k)=r-6\alpha_Rn_1-c_Rn_2,
 \qquad
 S_R(k)=s-\eta_Rn_1-\alpha_Rn_2.
 \label{eq:current-shifted-indices}
\end{equation}
Then
\begin{equation}
 P_R^{-\mathsf T}k=
 \left(n_1,
 \frac{m_R\mathcal R_R(k)+q_RS_R(k)}{D_R},n_2,
 \frac{-L_R\mathcal R_R(k)+6m_RS_R(k)}{D_R}\right).
 \label{cusp-spec-exact-dual-lattice-formula}
\end{equation}
The scalar Laplacian defined by the \(H^1\) Dirichlet form has operator
domain \(H^2(F_R)\), compact resolvent, and eigenvalues
\begin{equation}
 \lambda_R(k)=4\pi^2\lVert P_R^{-\mathsf T}k\rVert^2,
 \qquad k\in\mathbb Z^4,
 \label{eq:current-scalar-spectrum}
\end{equation}
with Fourier multiplicity.

Finally,
\begin{equation}
 D_R=R^2+(2H-B)R+H(H-B)+6m^2+O(Re^{-2\pi R}).
 \label{cusp-spec-D-expansion}
\end{equation}
\end{proposition}

\begin{proof}
The identity
\[
 Z=s\begin{pmatrix}0&1\\-1&0\end{pmatrix}
 +\begin{pmatrix}6\mu&h\\b-h&\mu\end{pmatrix},
\]
which follows from \eqref{eq:cusp-s-coordinate} and \(b=\beta+\tau\),
gives \eqref{cusp-spec-real-period-matrix} by taking real and imaginary
parts in the fixed coordinate order.  Expansion along its last two columns
gives
\[
 \det P_R=-L_Rq_R-6m_R^2=-D_R.
\]
The sign \(D_R>0\) is exactly the earlier condition \(D<0\), since
\(D=\operatorname{Im}\beta_R-6m_R^2/L_R=-D_R/L_R\) and \(L_R>0\).

Multiplying the displayed candidate inverse transpose by \(P_R^{\mathsf T}\)
gives the identity matrix entry by entry: the two nontrivial \(2\times2\)
minors reduce to \(L_Rq_R+6m_R^2=D_R\), and the real-part terms cancel by
the definitions of \(c_R\) and \(S_R(k)\).  This proves
\eqref{cusp-spec-full-inverse-transpose}, and multiplication by \(k\)
gives \eqref{cusp-spec-exact-dual-lattice-formula}.

The characters
\(x\mapsto\exp(2\pi i\langle P_R^{-\mathsf T}k,x\rangle)\) form the complete
Fourier basis of the compact quotient.  Applying the nonnegative Euclidean
Laplacian gives \eqref{eq:current-scalar-spectrum}; closed-form theory and
elliptic regularity give the stated form and operator domains, and Rellich
compactness gives compact resolvent.  Finally, substituting
\(L_R=R+H_R\), \(q_R=R+H_R-B_R\), and the three holomorphic expansions into
\(D_R=L_Rq_R+6m_R^2\) gives the last formula with every constant and sign
retained.
\end{proof}
